\textbf{Задача 4.}  Не опираясь на полиномиальный алгоритм проверки простоты, докажите, что множество простых чисел лежит в $\np \cap \conp$.

\begin{proof} Очевидно, что множество простых чисел лежит в $\mathbf{coNP}$, так как в качестве сертификата можно предъявлять разложение составного числа в произведение двух его делителей. Покажем, что проверка числа на простоту лежит в $\mathbf{NP}$. Заметим, что число $N$ простое $\lra \Z_{N}^*$ имеет порядок $N - 1$. Причем в этом случае $\Z_{N}^*$ будет циклической. Будем предъявлять в качестве сертификата первообразный корень $a$ и разложение $N - 1 = p_{1}p_{2} \dots p_{k}$ в произведение простых делителей. Для проверки простоты $N$ нужно проверить корректность разложения $N - 1$ и $\forall i$ простоту $p_i$, а также то, что $a^{(N - 1) / p_i} \not\equiv 1 \pmod{p}$ (это является критерием того, что $a$ --- первообразный корень). Первое и последнее условия проверяются за полином от $\log{N}$. Чтобы проверить условие простоты для $p_i$, воспользуемся рекурсией. А именно, добавим в сертификат сертификаты для простых делителей. Таким образом, сертификат $S(N)$ для $N$ будет выглядеть как $S(N) = (a, p_1, \dots, p_k, S(p_1), \dots, S(p_k))$. Простоту делителей будем проверять рекурсивно. Докажем индукцией, что в этой рекурсии будет не более $C \log{N} - C$ вызовов, где $C$ выбрано достаточно большим, чтобы утверждение было верно для маленьких $N$. Далее совершим переход, предположив, что для $p_i$ утверждение верно: количество вызовов для $N$ не больше
\[
1 + C \sum_{i=1}^k (\log{p_i} - 1) = 1 + C \log{(N - 1)} - Ck \le C\log{N} - C,
\]
так как $C$ можно изначально взять таким, чтобы $1 - Ck \le 1 - 2C \le -C$ ($k \ge 2$).

Получается, всего вызовов порядка логарифм от входа, при этом каждому рекурсивному вызову соответствует первообразный корень $a$ и набор простых чисел в записи сертификата, и всё это требует $\log{N}$ бит для записи. Значит, длина сертификата полиномиальна, и рекурсивный вызов требует лишь полиномиального времени. Таким образом, $\mathsf{PRIMES} \in \mathbf{NP}$.

\end{proof}